\section{Related Work}
\label{sec:related}

\paragraph{Global optimization and sequencing of chiller plants.}
Plant-level optimization for chiller systems has long targeted equipment sequencing, setpoint selection, and economic dispatch. Niu et al.~\citep{niu2023global} showed that even under existing factory operating constraints, a model-based global optimization workflow can achieve substantial savings without interfering with production. Bai et al.~\citep{bai2024modelica} further highlighted the interaction among sequencing, load allocation, and parameter optimization in heterogeneous multi-chiller plants. These studies provide strong evidence that system-level coordination matters, but they are primarily driven by simulation-centric optimization loops and do not directly emphasize an open, compact MILP implementation for robust day-ahead scheduling.

\paragraph{Modelica-based calibration and digital twins.}
Equation-based and object-oriented modeling frameworks have become important for diagnosing performance degradation and evaluating candidate operating policies. Fu et al.~\citep{fu2019equation} demonstrated how a calibrated Modelica model can identify hidden inefficiencies and quantify energy-saving opportunities in data-center cooling systems. Such models are valuable because they support high-fidelity representation and plant-specific interpretation. At the same time, the resulting optimization workflow can be tightly coupled to a particular simulator and calibration effort, which may limit reproducibility when the research goal is methodological scheduling rather than a site-specific digital twin.

\paragraph{Predictive control and storage-coupled cooling.}
Another major line of work develops model predictive control (MPC) or learning-enhanced closed-loop control for cooling systems. In data-center applications, Zhu et al.~\citep{zhu2023storage} proposed an MPC framework for chillers coupled with cold-water storage and reported meaningful annual cost and PUE benefits. Zhao et al.~\citep{zhao2024mpc} considered whole-system energy conservation in multi-chiller data-center control by combining predictive models with particle swarm optimization. These approaches are compelling for online operation, but they usually require forecasting stacks, rolling re-optimization, and tighter integration between plant dynamics and the controller. Our work instead targets the planning layer and asks how much value can be obtained from a simpler but still system-aware robust scheduling formulation.

\paragraph{Robust and stochastic chiller scheduling.}
Uncertainty-aware formulations recognize that cooling demand and tariff conditions are imperfectly known in advance. Saeedi et al.~\citep{saeedi2019robust} and Sadat-Mohammadi et al.~\citep{sadatmohammadi2020robust} showed that robust optimization can reduce sensitivity to uncertain demand and pricing in multi-chiller settings, including systems with chilled-water storage. These papers are closest in spirit to our study. The gap we address is the combination of three ingredients in a single reproducible research scaffold: heterogeneous chiller commitment, explicit auxiliary-equipment coupling, and a benchmark-oriented Pyomo implementation that can be extended toward journal-quality ablation and visualization workflows.

Overall, the existing literature supports the importance of system-level optimization, but it still leaves room for transparent and extensible scheduling frameworks that sit between heuristic operation and simulator-heavy predictive control. This paper is intended to fill that gap for the day-ahead planning problem.
